\documentclass[12pt]{article}

\usepackage[letterpaper,margin=1in]{geometry}
\usepackage{amsmath, amssymb, graphics, setspace}
\usepackage{graphicx} 
\graphicspath{{./figs/}}
\usepackage{tcolorbox}
\usepackage{bm}
\usepackage[framed,numbered,autolinebreaks,useliterate]{mcode}
\usepackage{hyperref}
\hypersetup{
    colorlinks=true,
    linkcolor=blue,
    filecolor=magenta,      
    urlcolor=cyan,
    pdftitle={Homework 13 Solutions},
    pdfpagemode=FullScreen,
}

\def\mathbi#1{\textbf{\em #1}}

\newcommand*{\I}{\imath} % imaginary unit i

\begin{document}
\doublespacing	

\title{CE 401/5501 Homework 13 Solutions\\
\vspace{2in}
\textbf{Name: \rule{2in}{0.15mm}}}
\author{}
\maketitle

\newpage
\doublespacing

\section*{Q1}
Please refer to Lecture 16 notes for context.
\begin{enumerate}
  \item \textbf{Era of Expert Systems:} Expert systems emerged as a major AI research topic in the late 1970s and throughout the 1980s.  This period, often called the 
  "Knowledge Engineering Era," was characterized by significant industry investment and commercial success of rule-based systems in domains such as medical diagnosis (e.g., MYCIN) and chemical process control (e.g., DENDRAL).

  \item \textbf{Methodological Components:}
    \begin{itemize}
      \item \emph{Knowledge Base:} A curated set of if--then production rules encoding domain expertise.
      \item \emph{Inference Engine:} Mechanisms for forward and/or backward chaining to derive conclusions from the knowledge base.
      \item \emph{Knowledge Acquisition:} Processes (often interviews and manual rule authoring) for capturing expert knowledge.
      \item \emph{Working Memory:} A database of facts and intermediate conclusions maintained during inference.
      \item \emph{Explanation Facility:} Modules that generate human-readable justifications for the system’s reasoning.
      \item \emph{User Interface:} Interactive shells allowing users to query the system and inspect reasoning steps.
    \end{itemize}

  \item \textbf{Challenges and Limitations:}
    \begin{itemize}
      \item \emph{Knowledge Acquisition Bottleneck:} Eliciting and formalizing expert rules was time-consuming and expensive.
      \item \emph{Brittleness:} Systems failed catastrophically outside their narrow rule coverage.
      \item \emph{Maintenance Overhead:} Rule bases grew large and became difficult to debug or update.
      \item \emph{Uncertainty Handling:} Early systems lacked robust methods to manage imprecise or probabilistic information.
      \item \emph{Scalability:} The rule explosion problem limited applicability to large, complex domains.
    \end{itemize}

  \item \textbf{Strengths and Weaknesses:}
    \begin{itemize}
      \item \emph{Classic AI (Expert Systems):}
        \begin{itemize}
          \item \textbf{Strengths:} Transparent reasoning, explicit domain knowledge, reliable performance within a defined scope, and explainability.
          \item \textbf{Weaknesses:} Inflexible, high development and maintenance cost, poor handling of uncertainty, and limited learning capability.
        \end{itemize}
      \item \emph{Modern Generative AI (LLMs):}
        \begin{itemize}
          \item \textbf{Strengths:} Broad coverage of language tasks, few-/zero-shot learning, adaptability, fluent natural-language generation, and scalability via pretraining.
          \item \textbf{Weaknesses:} Opaque reasoning ("black box"), potential for hallucinations, difficulty in precise control, large compute and data requirements, and limited inherent explainability.
        \end{itemize}
    \end{itemize}

\end{enumerate}

\newpage
\section*{Q2}
Please refer to Lecture 17 notes for context.
\begin{enumerate}
  \item \textbf{Transformers and the Move to Zero-/Few-Shot Frameworks:} 
    The transformer architecture introduced by Vaswani \emph{et al.} (2017) powered the shift.  Key innovations included:
    \begin{itemize}
      \item \emph{Self-Attention Mechanism:} Allowing each token to attend to all others, capturing long-range dependencies efficiently.
      \item \emph{Multi-Head Attention:} Learning diverse relational representations in parallel attention heads.
      \item \emph{Positional Encoding:} Injecting token order information without recurrent structures.
      \item \emph{Layer Normalization and Residual Connections:} Stabilizing training of very deep networks.
      \item \emph{Parallelizable Computation:} Enabling large-batch pretraining on massive text corpora, which underlies few-shot transfer ability.
    \end{itemize}

  \item \textbf{Basic Components of a Transformer-Based LLM:}
    \begin{enumerate}
      \item \emph{Input Embedding Layer:} Maps tokens to continuous vectors.
      \item \emph{Positional Encodings:} Adds sequence-order information.
      \item \emph{Stack of Transformer Blocks:}
        \begin{itemize}
          \item Multi-Head Self-Attention sublayer.
          \item Positionwise Feed-Forward Network.
          \item Layer Normalization and Residual Connections around each sublayer.
        \end{itemize}
      \item \emph{Output Projection \& Softmax:} Projects final hidden states to vocabulary logits and applies softmax for token prediction.
      \item \emph{Learned Parameters:} Query, Key, Value matrices for each head; feed-forward weights; embeddings; layer norms; etc.
    \end{enumerate}

  \item \textbf{Modern LLM Embedding Dimensionality:} 
    Typical embedding dimensions range from 768 (e.g., BERT-base) up to 12,288 (e.g., GPT-3).  Larger dimensions provide greater representational capacity, capturing more semantic nuance at the cost of increased computation and memory.

  \item \textbf{Chain-of-Thought Prompt for Cubic Maximization:}
    \begin{tcolorbox}[colback=gray!5,colframe=gray!40]
    "Let's solve the maximum/minimum of the cubic function
    \begin{equation*}
      y(x)=a_3 x^3 + a_2 x^2 + a_1 x + a_0.
    \end{equation*}
    \begin{enumerate}
      \item Recall that extrema occur when the first derivative is zero. Compute:
        \begin{equation*}
          \frac{dy}{dx}=3a_3 x^2 + 2a_2 x + a_1.
        \end{equation*}
      \item Set the derivative equal to zero and solve the quadratic:
        \begin{equation*}
          3a_3 x^2 + 2a_2 x + a_1 = 0 \quad \Rightarrow \quad x=\frac{-2a_2 \pm \sqrt{4a_2^2 -12a_3 a_1}}{6a_3}.
        \end{equation*}
      \item Compute the second derivative:
        \begin{equation*}
          \frac{d^2y}{dx^2}=6a_3 x + 2a_2.
        \end{equation*}
      \item Evaluate $d^2y/dx^2$ at each critical point to classify as max (negative) or min (positive).
      \item Finally, plug the critical $x$ values back into $y(x)$ and compare to determine the global maximum.
    \end{enumerate}"
    \end{tcolorbox}

\end{enumerate}

\end{document}
